\begin{table}[!htbp]
\centering
\caption{Weekend social-media use and paired fixed-effects--DML comparisons}
\label{tab:robustness_weekend}
\small
\setlength{\tabcolsep}{5pt}
\begin{tabular*}{\textwidth}{@{\extracolsep{\fill}}lccccc}
\hline
Statistic & Fixed effects & Elastic net & Lasso & Random forest & Neural net
\\
\hline
\multicolumn{6}{l}{\textit{Panel A: Ordered-treatment FE versus PLR-DML}} \\
Weekend-use estimate & -0.011 & 0.025 & 0.025 & 0.023 & 0.032
\\
 & (0.034) & (0.024) & (0.024) & (0.024) & (0.023)
\\
$t$-statistic & -- & -1.038 & -1.026 & -0.953 & -1.214
\\
$p$-value & -- & 0.299 & 0.305 & 0.341 & 0.225
\\[0.5ex]
\multicolumn{6}{l}{\textit{Panel B: Binary-treatment FE versus IRM-DML}} \\
High weekend use estimate & 0.060 & 0.078* & 0.073 & 0.082* & --
\\
 & (0.057) & (0.046) & (0.045) & (0.047) & --
\\
$t$-statistic & -- & -0.283 & -0.201 & -0.333 & --
\\
$p$-value & -- & 0.777 & 0.841 & 0.739 & --
\\
\hline
\end{tabular*}
\begin{minipage}{0.97\textwidth}
\footnotesize Notes: Standard errors are reported in parentheses below the coefficients. Panel A uses the ordered five-category weekend-use measure. Panel B defines high weekend use as at least four hours and reports the IRM average treatment effect. The fixed-effects models include respondent and wave effects; DML uses five-fold respondent-level cross-fitting. The neural-net IRM is not reported. Panel A uses 500 full respondent-cluster bootstrap re-estimations. Panel B uses 500 paired respondent-cluster influence-function bootstrap draws. The complete-case sample contains $N=4,331$ person-wave observations. $^{***}p<0.01$, $^{**}p<0.05$, $^{*}p<0.10$.
\end{minipage}
\end{table}
